\section{A complete constructive subcase: a finite Runge certificate}
\label{sec:runge-certificate}

The following result is an actual, uniform arithmetic theorem.  Its input is
a finite algebraic certificate, its output is an explicit box containing every
integral point, and its proof uses only integer arithmetic and bounded sums.
It supplies a substantial effective subcase of the requested theorem, including
smooth curves of arbitrarily large genus.  It does not supply the missing
certificate for an arbitrary curve to which Siegel's theorem applies.

Throughout this section a polynomial is a finite coefficient list, in ascending
order.  Trailing zero coefficients are removed; the empty list represents zero.
Its degree, coefficients, sum of absolute coefficients, products, and value at
an integer are primitive recursive operations.  Polynomial equalities below
are therefore finite, decidable coefficient identities.  Integers can be
represented by any fixed primitive recursive signed coding.  Statements using
absolute values are abbreviations for statements about the natural-number
absolute-value function on those codes.

\subsection{Two elementary bounds}

\begin{lemma}[An integer root bound]
\label{lem:runge-root-bound}
Let $G(X)=\sum_{i=0}^{r}g_iX^i$ be a nonzero integer polynomial.  If $r=0$,
it has no root.  If $r\geq1$ and $G(x)=0$ for an integer $x$, then
\[
 |x|\leq \sum_{i<r}|g_i|.
\]
\end{lemma}
\begin{proof}
The assertion is immediate for $x=0$.  Otherwise put $t=|x|\geq1$ and
$L=\sum_{i<r}|g_i|$.  Since the nonzero integer $g_r$ has $|g_r|\geq1$,
the equation $g_rx^r=-\sum_{i<r}g_ix^i$ gives
\[
 t^r\leq |g_r|t^r
   \leq\sum_{i<r}|g_i|t^i
   \leq L t^{r-1}.
\]
Cancellation of the positive integer $t^{r-1}$ gives $t\leq L$.  The constant
case says simply that a nonzero integer is not zero.  All sums are finite;
the comparison $t^i\leq t^{r-1}$ is proved by induction on the exponent.
\end{proof}

\begin{lemma}[A nonzero product]
\label{lem:runge-product-bound}
If $u,v$ are integers with $uv\neq0$, then
\[
 |uv|\geq\max\{|u|,|v|\},\qquad
 |u|+|v|\leq |uv|+1.
\]
Moreover, if $v-u=2z$, then $\max\{|u|,|v|\}\geq |z|$.
\end{lemma}
\begin{proof}
Write $U=|u|$ and $V=|v|$.  Both are positive integers, so $UV\geq U,V$.
Expanding $(U-1)(V-1)\geq0$ gives $U+V\leq UV+1$.
Finally,
$2|z|=|v-u|\leq U+V\leq2\max(U,V)$, and cancellation of $2$ proves the
last assertion.  These are finite inequalities, with no choice of a prime
factor and no factorization theorem needed.
\end{proof}

\subsection{The certificate theorem and its complete proof}

\begin{theorem}[Elementary effective Runge certificate]
\label{thm:runge-certificate}
Let $f,A,R\in\mathbb Z[X]$ and $D\geq1$ satisfy the coefficient identity
\[
                 D^2 f=A^2+R.
\tag{RC}
\]
Assume that $m=\deg A\geq1$, that $R\neq0$, and that $\deg R<m$.
Write
\[
 \alpha=\sum_{i<m}|a_i|,\qquad
 \beta=\sum_i|r_i|,\qquad H=\alpha+\beta,
 \qquad K=\sum_i|f_i|H^i.
\]
Then every integer solution of $y^2=f(x)$ satisfies
\[
                         |x|\leq H,
                    \qquad |y|\leq K.
\]
Consequently the integral points have a primitive recursive complete
enumeration, uniformly in $f,D,A,R$ satisfying the displayed conditions.
\end{theorem}

\begin{proof}
Because $R$ is nonzero, $\beta\geq1$, hence $H\geq1$.
Let $y^2=f(x)$ and put $t=|x|$.  Integer order is decidable.  It suffices to
show that its alternative $t>H$ leads to a contradiction.

First, $R(x)\neq0$.  This is immediate if $R$ is constant.  If its degree is
$r\geq1$, Lemma~\ref{lem:runge-root-bound} would otherwise give
$t\leq\sum_{i<r}|r_i|\leq\beta\leq H$.

Second, $A(x)$ is larger in absolute value than $R(x)$.  Since $t\geq1$ and
the leading coefficient of $A$ is a nonzero integer, the reverse triangle
inequality gives
\begin{align*}
 |A(x)|
 &\geq |a_m|t^m-\sum_{i<m}|a_i|t^i\\
 &\geq t^m-\alpha t^{m-1}
   =(t-\alpha)t^{m-1}
   >\beta t^{m-1}.
\end{align*}
The strict inequality follows from $t>\alpha+\beta$ and
$t^{m-1}\geq1$.  Every exponent occurring in $R$ is at most $m-1$, so
\[
                  |R(x)|\leq\beta t^{m-1}<|A(x)|.
\]

Now use (RC) at the particular integer $x$ and substitute $y^2=f(x)$:
\[
              (Dy-A(x))(Dy+A(x))=R(x).
\]
The two factors are nonzero integers.  Their difference is $2A(x)$, so
Lemma~\ref{lem:runge-product-bound} gives
\[
 |R(x)|\geq
 \max\{|Dy-A(x)|,|Dy+A(x)|\}\geq |A(x)|.
\]
This contradicts the preceding strict inequality.  We have proved $t\leq H$.

For the second coordinate,
\[
 y^2=f(x)\leq |f(x)|\leq\sum_i|f_i|\,|x|^i
                         \leq\sum_i|f_i|H^i=K.
\]
For every integer $y$, the natural number $n=|y|$ satisfies $n\leq n^2$:
separate $n=0$ from $n\geq1$ and multiply $1\leq n$ by $n$.
Thus $|y|\leq y^2\leq K$.

Finally, run two bounded loops through $-H\leq x\leq H$ and
$-K\leq y\leq K$, retaining exactly the pairs for which the decidable
identity $y^2=f(x)$ holds.  Every retained pair is a solution by the test.
Every solution is retained by the two bounds just proved.  This produces a
finite list with both soundness and completeness, rather than merely a
numerical bound on the number of possible points.
\end{proof}

The use of a contradiction in this proof is intuitionistically harmless.
We did not infer an existential witness from the negation of its negation.
We computed $H$ and $K$ before considering a point, and used the decidable
integer dichotomy $t\leq H$ or $t>H$ to verify their universal correctness.
Similarly, the conclusion $R(x)\neq0$ follows by excluding equality in a
decidable integer calculation.  No principle for choosing from an infinite
family of finite sets occurs.

The factorization step explains the scope of the method.  The certificate
forces $Dy$ to be close to either $A(x)$ or $-A(x)$ at large $|x|$.  The two
integer factors cannot both remain nonzero while their product has lower
growth than their difference.  An arbitrary affine curve need not admit
this particular description; the theorem contains no algorithm for
manufacturing it in that generality.

The domain of the coordinates is equally essential.  The proof uses the
fact that the absolute value of a nonzero \emph{integer} is at least one.
For rational factors this fails: $u$ may be an arbitrarily small positive
fraction while $v$ is large.  Thus the same displayed polynomial identity
does not prove finiteness of rational points.  Nor does it automatically
prove an analogous assertion for coordinates whose denominators are
supported on a prescribed set of primes.  In that setting the product
argument would have to control several absolute values simultaneously;
the ordinary real absolute value alone no longer gives the required
lower bound for each factor.  The theorem here is deliberately stated
over $\mathbb Z$ with its fixed affine coordinates.

This distinction also explains why a birational change of variables must
not be inserted silently.  A rational transformation can introduce
denominators and can omit exceptional points.  To transfer the explicit
box to a different equation, one would have to give the transformation,
prove which integral points it covers, bound any denominators it creates,
and handle all exceptional fibers.  None of those obligations is needed
for the displayed families, whose integral points are bounded directly in
their original coordinates.  A polynomial identity checked on the actual
input equation is therefore part of the certificate, rather than an
informal assertion that the curve resembles an already-solved example.

\subsection{Constructing the certificate by completing a polynomial square}

The certificate is not an oracle for the following broad family.  Suppose
$f$ has degree $2m\geq2$ and its leading coefficient is a positive integer
square $s^2$, where $s\geq1$.  Construct a rational polynomial
\[
                      Q(X)=sX^m+\sum_{j<m}q_jX^j
\]
so that all coefficients of $f-Q^2$ in degrees $m,m+1,\ldots,2m$ vanish.
Start with $q_m=s$.  For $j=m-1,m-2,\ldots,0$, set
\[
 q_j=\frac{1}{2s}
 \left(f_{m+j}-
   \sum_{\substack{j<i<m,\ j<k<m\\i+k=m+j}}q_iq_k\right).
\tag{SQ}
\]
The sum uses only coefficients already defined.  Indeed, a term of degree
$m+j$ in $Q^2$ either uses one coefficient of degree $m$ and the coefficient
of degree $j$, or has both indices strictly greater than $j$.  The former
two terms total $2s q_j$.  Thus (SQ) sets the coefficient in degree $m+j$
equal to $f_{m+j}$ and leaves every previously matched higher coefficient
unchanged.  A finite downward induction verifies
\[
                         \deg(f-Q^2)<m
\]
unless $f-Q^2=0$, when the same coefficient assertion holds.

Take any positive common multiple $D$ of the denominators of the $q_j$.
Their product is sufficient; taking their least common multiple improves
the numerical size.  Define
\[
                        A=DQ,\qquad R=D^2f-A^2.
\]
These are integer polynomials and satisfy (RC).  If $R\neq0$, its degree is
less than $m$, so Theorem~\ref{thm:runge-certificate} applies immediately.
If $R=0$, the constructor reports the exact exceptional case
$f=Q^2$ in $\mathbb Q[X]$ and makes no finiteness claim.  Conversely, if
$f$ is a rational polynomial square, its square root with leading
coefficient $s$ has to obey the same triangular recursion, so this test
does detect that case.  In particular every nonsquare monic integer
polynomial of positive even degree is covered.  Every squarefree such
polynomial is nonsquare and therefore covered as well.

This construction uses exactly $m$ rational coefficient steps.  Each step
uses a bounded sum and rational addition, multiplication, and division by
the explicitly nonzero integer $2s$.  Denominators can be represented
without reduction, or reduced by the primitive recursive Euclidean
algorithm.  Clearing denominators, checking that the remainder is nonzero,
and checking the final coefficient identity are finite arithmetic
operations.  There is no search for a height bound with an unknown stopping
criterion.  The bounds can be large, and no claim of practical efficiency
is part of the theorem.

The restrictions should be kept visible.  For example, the displayed
constructor does not handle an odd-degree cubic or an even-degree
polynomial with a nonsquare positive leading coefficient.  Such a
rejection says that this constructor does not provide the certificate;
it says nothing about whether another method proves finiteness for the
curve.  Conversely, a zero remainder must not be ignored: the integer
polynomial $f(X)=(X^2+1)^2$ has the infinitely many points
$(x,\pm(x^2+1))$.

\subsection{A complete family in every positive hyperelliptic genus}

\begin{theorem}[Uniform exact classification]
\label{thm:runge-all-genera}
For every integer $m\geq2$, the integral solutions of
\[
                         y^2=x^{2m}+x+1
\]
are exactly
\[
                    (-1,-1),\quad(-1,1),\quad(0,-1),\quad(0,1).
\]
\end{theorem}

\begin{proof}
Use $D=1$, $A=X^m$, and $R=X+1$.  Then
$D^2f=A^2+R$, $\deg R=1<m$, $\alpha=0$, and $\beta=2$.
The certificate theorem gives $|x|\leq2$.  It remains to check just five
integers.

For $x=-1$ or $x=0$, the right side is $1$, giving exactly $y=\pm1$.
For $x=1$ it is $3$, strictly between the consecutive squares $1$ and $4$.
For $x=2$, put $k=2^m\geq4$.  The right side is $k^2+3$, and
\[
                         k^2<k^2+3<(k+1)^2,
\]
because $(k+1)^2-k^2=2k+1\geq9$.  For $x=-2$ it is $k^2-1$, and
\[
                        (k-1)^2<k^2-1<k^2,
\]
because $k^2-1-(k-1)^2=2k-2>0$.
An integer square cannot lie strictly between consecutive natural-number
squares: its absolute square root would have to be a natural number
strictly between the two consecutive integers.  This eliminates the
remaining values of $x$.  Substitution checks each of the four listed
points, completing both directions.
\end{proof}

Here the constant list of four points is already a uniform constructive
output; no growth of the list with the genus parameter is necessary.
The exponent $2m$ is part of the input, and exponentiation is primitive
recursive.  The argument uses only the finite coefficient certificate and
inequalities between integers.  It does not assume a general theorem
about rational points or integral points on high-genus curves.

For completeness, these examples really do represent smooth positive-genus
curves, rather than a collection of singular rational equations.  Set
$f_m=X^{2m}+X+1$.  A common zero $\xi$ of $f_m$ and its derivative would
satisfy the polynomial identity
\[
       2m f_m(X)-X f_m'(X)=(2m-1)X+2m
\]
and hence $\xi=-2m/(2m-1)$.  This is a reduced nonintegral rational number
for $m\geq2$.  A rational root of a monic integer polynomial is an integer:
writing it as $p/q$ in lowest terms and clearing denominators shows
$q\mid p^{\deg f}$, so $q=1$.  Thus no such common zero exists, and $f_m$
is squarefree.  An alternative finite certificate is reduction modulo
$2$, where $f_m'$ becomes the constant polynomial $1$.

The affine equation is nonsingular in characteristic zero: a singular
point would have $2y=0$, $f_m(x)=0$, and $f_m'(x)=0$, which has just been
excluded.  At infinity put $t=1/x$ and $w=y/x^m$.  The equation becomes
\[
                         w^2=1+t^{2m-1}+t^{2m}.
\]
At $t=0$ there are exactly the two rational points $w=1$ and $w=-1$;
both are smooth because the derivative with respect to $w$ is $2w\neq0$.
The smooth completion is the double cover of the projective line branched
at the $2m$ simple finite roots of $f_m$, with no ramification at infinity.
The usual hyperelliptic genus calculation is
$2g-2=2(-2)+2m$, giving $g=m-1$.  This geometric identification explains
the range of the examples.  The arithmetic classification above is
independent of that genus calculation and requires none of its geometry
as an additional axiom of the arithmetic proof.

\subsection{A sharper quartic calculation and a divisor algorithm}

The symmetric quartic family admits an especially transparent certificate
and sharper numerical bounds.  Let $a,b\in\mathbb Z$ and put
$\Delta=a^2-4b$.  Assume $\Delta\neq0$.  A solution of
\[
                         y^2=x^4+ax^2+b
\]
satisfies
\[
 (2x^2+a-2y)(2x^2+a+2y)=\Delta.
\tag{QF}
\]
Write the two factors as $u,v$ and let $d=|\Delta|$.
Both are nonzero and $|u|+|v|\leq d+1$ by
Lemma~\ref{lem:runge-product-bound}.  Since
$u+v=4x^2+2a$ and $v-u=4y$, we obtain the useful signed bounds
\[
                    4x^2\leq d+1-2a,
                 \qquad 4|y|\leq d+1.
\tag{QB}
\]
If the first right side is negative, the equation has no integral points.
Otherwise one may take
\[
 H=\left\lfloor\sqrt{\left\lfloor(d+1-2a)/4\right\rfloor}\right\rfloor,
 \qquad K=\left\lfloor(d+1)/4\right\rfloor.
\]
For the more general certificate theorem one could simply take
$D=2$, $A=2X^2+a$, and $R=4b-a^2$, giving the less sharp but still explicit
bound $|x|\leq |a|+|\Delta|$.

There is also a complete algorithm that avoids scanning this rectangle.
Enumerate the nonzero integers $u$ in $[-d,d]$ and retain those dividing
$\Delta$.  Set $v=\Delta/u$.  For each ordered factor pair test whether
\[
           z=\frac{u+v-2a}{4},\qquad y=\frac{v-u}{4}
\]
are integers, whether $z\geq0$, and whether $z$ is an integer square.
If $z=x_0^2$, return $(x_0,y)$ and $(-x_0,y)$, removing duplicates.
Every solution supplies exactly such a factor pair by (QF), so none is
missed.  Conversely the displayed reconstruction gives
$2z+a=(u+v)/2$ and $2y=(v-u)/2$; their difference of squares is
$uv=\Delta$.  Expanding and substituting $\Delta=a^2-4b$ yields
$y^2=z^2+az+b$, the original equation with $z=x_0^2$.  Thus every returned
pair is valid.  Divisibility, the congruences modulo $4$, and the square
test are all bounded decisions.

For a concrete positive-genus example, choose $a=7$ and $b=1$, so
$\Delta=45$.  The bounds give $x^2\leq8$ and $|y|\leq11$.
In (QF), $u+v=4x^2+14>0$ and $uv=45>0$, so both factors are positive.
Up to order the positive factor pairs are
$(1,45)$, $(3,15)$, and $(5,9)$.  They give respectively
$z=8,1,0$.  The first is not a square; the latter two give exactly
\[
                         (0,\pm1),\qquad(\pm1,\pm3).
\]
This is a complete list of six points.  It is verified by the factor
argument, not by interpreting a finite experimental search as complete.

More generally the quartic polynomial $X^4+aX^2+b$ is squarefree when
$b\Delta\neq0$.  A repeated root would have either $x=0$, forcing $b=0$,
or $2x^2+a=0$, forcing $b-a^2/4=0$.  Both alternatives are excluded.
Its smooth completion is therefore a genus-one curve, with two rational
points at infinity.  The assumptions hold for $a=7,b=1$.
The arithmetic bound itself only requires $\Delta\neq0$, so it also
handles some singular examples.  When $\Delta=0$, integrality forces $a$
to be even and
$X^4+aX^2+b=(X^2+a/2)^2$; in that case there are infinitely many integral
points.  The excluded zero remainder is therefore a necessary distinction
for this particular family.

